\section*{Discussion}

\subsection*{Haplotype match-length as a principled extension of existing introgression frameworks}

ArchaicPainter (\AP{}) demonstrates that the Li \& Stephens haplotype copying
framework, long established for modern local ancestry inference, extends naturally
to the archaic introgression detection problem.
The central theoretical argument --- that an 11$\times$ coalescent tract-length
contrast between introgressed and non-introgressed segments should translate into
substantial detection gains --- is borne out empirically: \AP{} achieves F1 $=
0.480 \pm 0.095$ compared to $0.009 \pm 0.003$ for a Poisson density baseline
($56\times$ improvement), directly attributable to the match-length signal
captured in the HMM transition matrix.
This gap is not merely a numerical result; it suggests that archaic introgression
methods based primarily on density-counting principles may operate with a
significant power deficit by discarding the linkage disequilibrium structure
that encodes admixture history --- the same structure that has driven
modern local ancestry inference for two decades.

The $56\times$ F1 advantage should be interpreted carefully.
It is measured against a simplified Poisson density baseline that represents
the density-counting principle, not against HMMix or DAIseg directly; the latter
employ HMMs with their own transition structures and may capture some spatial
correlation implicitly.
Nevertheless, the density-counting principle is the foundational design choice
of all existing archaic introgression HMMs, and the scale of the F1 gap strongly
suggests that explicit match-length modeling deserves serious investigation
in the context of more sophisticated reference tools.
We regard the 50-replicate simulation benchmark as a controlled demonstration
of the principle, sufficient for publication at PLOS Genetics as a proof of concept,
and defer direct head-to-head F1 benchmarking against HMMix and DAIseg to
future work with access to those tools' code and standardized benchmarking
frameworks.

\subsection*{Emission model selection and the reference--donor divergence problem}

A recurring challenge in applying haplotype-based methods to ancient DNA analysis
is reference--donor divergence: the available archaic reference genome (e.g.,
Vindija Neanderthal 33.19) is an imperfect proxy for the true introgressing
population.
\AP{} addresses this via the positive-only emission mode, which treats mismatches
as uninformative rather than penalizing them.
The ablation result (F1 $= 0.000$ under positive-only on simulation, where the
reference exactly matches the donor) provides a precise characterization of when
each mode should be used: the standard bidirectional model is optimal whenever the
reference tracks the true donor, while the positive-only mode avoids false negatives
introduced when mismatches between the query and an imperfect reference are
incorrectly interpreted as evidence against introgression.

This emission mode switch is conceptually related to approaches in ancient DNA
analysis that handle post-mortem damage and reference bias through asymmetric
likelihood models~\citep{jonsson2013mapdamage}.
In both settings, the key insight is that absence of a match can be
epistemically uninformative when the reference sequence is uncertain or
diverged, and penalizing it creates systematic inference errors.
The distinction between contexts in which mismatches \emph{are} informative (simulation,
where the reference is the exact donor) versus when they are \emph{not} informative
(real data with a diverged reference) is fundamental to method deployment and
is not always made explicit in archaic introgression literature.
\AP{} makes this choice transparent and provides ablation validation of the
two regimes.

\subsection*{Per-haplotype resolution as a qualitative advance over diploid methods}

The comparison with IBDmix highlights a dimension of progress distinct from
raw sensitivity: \AP{} operates at the level of phased haplotypes rather than
diploid individuals.
This haplotype resolution enables analyses that diploid methods cannot support:
distinguishing which of the two chromosomes within an individual carries an
introgressed segment, coupling introgressed haplotype identity across individuals
for population-level haplotype network analysis, and in principle integrating
with modern phasing-based methods such as SHAPEIT~\citep{delaneau2019accurate}
for jointly phased ancient--modern data.
The 37 per-haplotype segments detected by \AP{} on chr21:10--20~Mb versus
345 diploid segments by IBDmix at the same LOD threshold are not simply two
points on the same receiver operating characteristic curve; they represent
fundamentally different units of biological measurement.
Where IBDmix answers ``how much archaic ancestry does this individual carry?'',
\AP{} answers ``which specific haplotype backbone in this individual is of
archaic origin, and what is its extent?''

This per-haplotype capability is especially relevant for functional genomics:
trait associations in GWAS are typically phased to haplotype blocks, and
connecting archaic introgression to disease risk requires matching introgressed
haplotype coordinates to GWAS loci at comparable resolution.
The 79 named genes detected in our chr21 analysis, including disease-associated
loci from immune, neurological, and metabolic pathways (Supplementary Table~S1),
demonstrate that per-haplotype coordinates from \AP{} can be directly input
to downstream GWAS locus intersection without the imprecision introduced by
diploid-averaging.

\subsection*{The interferon receptor cluster and adaptive retention of Neanderthal immunity haplotypes}

The unsupervised recovery of the complete chr21 interferon receptor gene cluster
(\textit{IFNAR1}, \textit{IFNAR2}, \textit{IFNGR2}, \textit{IL10RB}) from
\AP{}-predicted Neanderthal haplotypes is the most biologically significant
finding of the real-data analysis.
These four genes define the molecular docking interface for the type-I interferon
(IFN-$\alpha/\beta$), type-II interferon (IFN-$\gamma$), and anti-inflammatory
interleukin-10/type-III interferon response pathways --- the three major arms of
the innate antiviral immune response.
This cluster was independently identified as a Neanderthal introgression target in
European populations by Vernot \& Akey~\citeyear{vernot2014resurrecting}
and Sankararaman et al.~\citeyear{sankararaman2014genomic};
McCoy et al.~\citeyear{mccoy2017impacts} subsequently showed that
Neanderthal haplotypes at these loci are associated with altered
transcript levels in multiple immune cell types.
Our rediscovery of the same cluster, using a completely different computational
approach applied without any annotation information, provides independent
corroboration of this signal and illustrates the methodological principle that
\AP{}'s match-length emissions are specific enough to recover locus-level
adaptive introgression candidates in unsupervised analysis.

The co-detection of this cluster in CEU but not prominently in CHB
is consistent with the geography of the interferon locus signal:
European-specific Neanderthal haplotype effects at immune loci have been
reported independently~\citep{dannemann2017contribution},
potentially reflecting differential positive selection for antiviral immunity
in the dispersal environments encountered by European-ancestral populations.
The population-specific patterns for other gene overlaps (\textit{DYRK1A} in CEU;
\textit{RUNX1} and \textit{ADAMTS5} in CHB) are similarly consistent with published
population genomics of chr21 adaptive introgression~\citep{vernot2016excavating},
suggesting that \AP{} reproduces known geography of introgression without being
parametrized to do so.

\subsection*{The multi-source attribution challenge: current limitations and path forward}

The proof-of-concept three-state attribution results (NEA F1 $= 0.260$, DEN
F1 $= 0.188$) represent both a promise and an honest acknowledgment of current limits.
The promise is that the same HMM framework that accurately detects Neanderthal
ancestry in the two-state case extends to simultaneous multi-source attribution
without architectural changes; the attribution accuracy ($\sim$26\% for NEA,
$\sim$15\% for DEN) is above chance, confirming that the three-state model
extracts some source-discriminating information.

The limits are real and primarily biological.
The $\sim$700-kya Neanderthal--Denisovan divergence~\citep{prufer2014complete}
is much more recent than the $\sim$6-Mya Homo--Pan split that provides a dense
background of diagnostic markers for modern local ancestry inference;
the smaller number of archaic-lineage-specific variants means more ambiguous
emission evidence for source disambiguation.
Improving attribution accuracy will require one or more of:
(1) higher-quality phased Denisovan reference genomes
(the Altai Denisovan~\citep{meyer2012high} is the only available genome and is
from a single individual on a small island),
(2) leveraging reads from Denisovan-like populations (e.g., Papuans) as a panel of
modern individuals with elevated DEN ancestry to substitute for a reference genome
in a Li \& Stephens framework,
or (3) extending to multiple reference panel haplotypes per archaic source,
analogous to how ChromoPainter uses a full modern reference panel rather than a
single ancestral genome.
The near-linear scalability of \AP{} means that extending the reference panel
is computationally feasible; the bottleneck is data availability for DEN references.

\subsection*{Limitations and scope}

Several limitations of the current work merit explicit acknowledgment.

First, our real-data analysis is restricted to chromosome~21.
Chromosome~21 is the smallest human autosome and carries a relatively modest
number of archaic-informative sites (1,036 Vindija-matching SNPs in 38~Mb,
$\sim$27 SNPs/Mb), which is an order of magnitude below the simulation density
($\sim$320 SNPs/Mb).
This informative site sparsity constrains segment resolution and likely reduces
F1 relative to simulation; a full genome-wide analysis would provide more informative
sites per chromosome and allow assessment of \AP{} performance at the per-gene level.

Second, no direct quantitative benchmark against HMMix or DAIseg was performed.
The comparison to a Poisson density baseline is an isolating experiment for the
match-length signal, not a competitive evaluation.
Reviewers familiar with HMMix (which uses LD-aware transition parameters in a
two-state HMM) may note that our baseline is deliberately simplified; a fair
comparison with HMMix on identical test sets would require access to that tool's
code and a jointly designed evaluation framework with matched parameters.
We are transparent about this limitation and encourage future benchmarking studies
to include \AP{} as a third point in this comparison.

Third, the positive-only emission mode used for real data has not been calibrated
against a real-data ground truth.
The decision to use positive-only is theoretically motivated (reference--donor
divergence) and ablation-validated in simulation, but the optimal emission weight
for Vindija's specific divergence from the introgressing Neanderthal population
remains an open empirical question.

Fourth, the simulation demography uses a simplified single-pulse admixture model.
Real admixture likely involved multiple waves~\citep{vernot2016excavating},
background selection removing deleterious introgressed alleles, and population
structure within archaic lineages.
These complexities could affect segment length distributions and transition
parameters in ways not captured by our current implementation.

Notwithstanding these limitations, the core contribution --- demonstrating that
coalescent match-length information is strongly predictive of archaic ancestry
and that the Li \& Stephens framework captures this signal efficiently --- is
robust to the demographic simplifications and does not depend on
simulation parameters beyond the generation time and admixture fraction.

\subsection*{Future directions}

The most immediately impactful extension is a full-genome sweep across the
1000~Genomes Project and/or the Simons Genome Diversity Panel.
Chromosome~21 yields 79 named genes; a rough extrapolation to the full autosomes (79 genes on the smallest
autosome) suggests on the order of $\sim$2,000--4,000 named genes in
\AP{}-detected introgressed segments genome-wide, broadly comparable in
scale to the Sankararaman et al.~\citeyear{sankararaman2014genomic} catalog
but at per-haplotype rather than per-individual resolution.

Integration of SparsePainter's PBWT-based sparse haplotype matching
algorithm~\citep{yang2025sparsepainter} could bring \AP{} to biobank scale
($n \gtrsim 500{,}000$ haplotypes); whether genome-wide archaic-informative
site density ($\sim$27 SNPs/Mb on chr21) is sufficient for robust per-segment
inference at this scale requires empirical validation.
Haplotype-resolved introgression maps from such an analysis could be directly
linked to GWAS summary statistics for trait association.
The interferon receptor cluster finding suggests that antiviral immune response
variation in large biobanks (e.g., the UK Biobank) may in part reflect
Neanderthal introgression haplotype segregation --- a hypothesis now testable
at scale with \AP{}.
